
\documentclass{article}
\usepackage{colortbl}
\usepackage{makecell}
\usepackage{multirow}
\usepackage{supertabular}

\begin{document}

\newcounter{utterance}

\centering \large Interaction Transcript for game `matchit\_ascii', experiment `same\_grid', episode 2 with NEW\_Reflect\_Agent.
\vspace{24pt}

{ \footnotesize  \setcounter{utterance}{1}
\setlength{\tabcolsep}{0pt}
\begin{supertabular}{c@{$\;$}|p{.15\linewidth}@{}p{.15\linewidth}p{.15\linewidth}p{.15\linewidth}p{.15\linewidth}p{.15\linewidth}}
    \# & $\;$A & \multicolumn{4}{c}{Game Master} & $\;\:$B\\
    \hline

    \theutterance \stepcounter{utterance}  
    & & \multicolumn{4}{p{0.6\linewidth}}{
        \cellcolor[rgb]{0.9,0.9,0.9}{
            \makecell[{{p{\linewidth}}}]{
                \texttt{\tiny{[P1$\langle$GM]}}
                \texttt{You are participating in a collaborative guessing game.} \\
\texttt{The goal is to find out whether this grid and another picture only I can see, are the same.} \\
\texttt{Please describe your grid first. Then, I will provide my description and we can ask each other questions about the grids to figure out whether they are the same. Now start your short grid description with "DESCRIPTION:" followed by the description. Do not add anything else. This is your grid:} \\
\texttt{▢ ▢ ▢ ▢ ▢} \\
\texttt{▢ X X X X} \\
\texttt{▢ ▢ ▢ X X} \\
\texttt{▢ X X X X} \\
\texttt{▢ X X X X} \\
            }
        }
    }
    & \\ \\

    \theutterance \stepcounter{utterance}  
    & \multicolumn{4}{p{0.45\linewidth}}{
        \cellcolor[rgb]{0.95,0.95,0.95}{
            \makecell[{{p{\linewidth}}}]{
                \texttt{\tiny{[P1$|$internal]}}
                \texttt{Reflections from previous episodes:} \\
\texttt{{-}} \\
\texttt{ Failed. My description was incomplete and did not capture the exact shape patterns. I only described overall X clusters but missed the precise layout—specifically, the distinct “T” and “U” shapes formed by the blank spaces. Next time, I would systematically count rows and columns and note the exact pattern of blanks to give a more precise structural description.} \\
\texttt{{-}} \\
\texttt{ Failed. I described only the central row's pattern, missing the full grid symmetry. This made my description too vague for accurate comparison. Next time, I would explicitly note the full grid layout: all X's except the three squares in the central row.} \\
            }
        }
    }
    & & \\ \\

    \theutterance \stepcounter{utterance}  
    & \multicolumn{4}{p{0.45\linewidth}}{
        \cellcolor[rgb]{0.95,0.95,0.95}{
            \makecell[{{p{\linewidth}}}]{
                \texttt{\tiny{[P1$|$internal]}}
                \texttt{{[}attempt 1 – valid{]}} \\
\texttt{DESCRIPTION: A 5x5 grid where the first row is all blank (▢), the second row has X's in columns 2{-}5, the third row has blanks in columns 1{-}3 and X's in columns 4{-}5, the fourth row has X's in columns 1{-}5, and the fifth row has X's in columns 1{-}5. The pattern is symmetric across the vertical axis, with the central row (third row) having a "U" shape of blanks on the left and X's on the right.} \\
            }
        }
    }
    & & \\ \\

    \theutterance \stepcounter{utterance}  
    & \multicolumn{4}{p{0.6\linewidth}}{
        \cellcolor[rgb]{0.9,0.9,0.9}{
            \makecell[{{p{\linewidth}}}]{
                \texttt{\tiny{[P1$\rangle$GM]}}
                \texttt{DESCRIPTION: A 5x5 grid where the first row is all blank (▢), the second row has X's in columns 2{-}5, the third row has blanks in columns 1{-}3 and X's in columns 4{-}5, the fourth row has X's in columns 1{-}5, and the fifth row has X's in columns 1{-}5. The pattern is symmetric across the vertical axis, with the central row (third row) having a "U" shape of blanks on the left and X's on the right.} \\
            }
        }
    }
    & & \\ \\

    \theutterance \stepcounter{utterance}  
    & & & \multicolumn{2}{p{0.3\linewidth}}{
        \cellcolor[rgb]{0.9,0.9,0.9}{
            \makecell[{{p{\linewidth}}}]{
                \texttt{\tiny{[GM$|$GM]}}
                \texttt{continue} \\
            }
        }
    }
    & & \\ \\

    \theutterance \stepcounter{utterance}  
    & & \multicolumn{4}{p{0.6\linewidth}}{
        \cellcolor[rgb]{0.9,0.9,0.9}{
            \makecell[{{p{\linewidth}}}]{
                \texttt{\tiny{[GM$\rangle$P2]}}
                \texttt{You are participating in a collaborative guessing game.} \\
\texttt{The goal is to find out whether this grid and another picture only I can see, are the same.} \\
\texttt{Please describe your grid first. Then, I will provide my description and we can ask each other questions about the grids to figure out whether they are the same. Now start your short grid description with "DESCRIPTION:" followed by the description. Do not add anything else. This is your grid:} \\
\texttt{▢ ▢ ▢ ▢ ▢} \\
\texttt{▢ X X X X} \\
\texttt{▢ ▢ ▢ X X} \\
\texttt{▢ X X X X} \\
\texttt{▢ X X X X} \\
            }
        }
    }
    & \\ \\

    \theutterance \stepcounter{utterance}  
    & & & \multicolumn{4}{p{0.45\linewidth}}{
        \cellcolor[rgb]{0.95,0.95,0.95}{
            \makecell[{{p{\linewidth}}}]{
                \texttt{\tiny{[P2$|$internal]}}
                \texttt{Reflections from previous episodes:} \\
\texttt{{-}} \\
\texttt{ Failed. I gave a full description immediately, leaving no room for strategic questioning. This allowed the user to simply copy my description without revealing any independent details about their grid. Next time, I should give a partial or ambiguous description first to force the user to reveal their own grid's specifics through their questions and answers.} \\
\texttt{{-}} \\
\texttt{ FAIL. I gave a description but then the user described their own grid instead of asking a question, breaking the intended turn{-}taking. I should have explicitly asked a question after my description to keep the interaction on track. Next time, I will ensure I follow the instruction to ask a question immediately after providing my description.} \\
            }
        }
    }
     \\ \\

    \theutterance \stepcounter{utterance}  
    & & & \multicolumn{4}{p{0.45\linewidth}}{
        \cellcolor[rgb]{0.95,0.95,0.95}{
            \makecell[{{p{\linewidth}}}]{
                \texttt{\tiny{[P2$|$internal]}}
                \texttt{{[}attempt 1 – valid{]}} \\
\texttt{DESCRIPTION: A 5x5 grid with the following pattern:  } \\
\texttt{{-} Row 1: All empty (□ □ □ □ □)  } \\
\texttt{{-} Row 2: First cell empty, then four X's (□ X X X X)  } \\
\texttt{{-} Row 3: First three empty, last two X's (□ □ □ X X)  } \\
\texttt{{-} Row 4: All X's (X X X X X)  } \\
\texttt{{-} Row 5: All X's (X X X X X)  } \\
\\ 
\texttt{Now, please describe your grid.} \\
            }
        }
    }
     \\ \\

    \theutterance \stepcounter{utterance}  
    & & & \multicolumn{4}{p{0.6\linewidth}}{
        \cellcolor[rgb]{0.9,0.9,0.9}{
            \makecell[{{p{\linewidth}}}]{
                \texttt{\tiny{[GM$\langle$P2]}}
                \texttt{DESCRIPTION: A 5x5 grid with the following pattern:  } \\
\texttt{{-} Row 1: All empty (□ □ □ □ □)  } \\
\texttt{{-} Row 2: First cell empty, then four X's (□ X X X X)  } \\
\texttt{{-} Row 3: First three empty, last two X's (□ □ □ X X)  } \\
\texttt{{-} Row 4: All X's (X X X X X)  } \\
\texttt{{-} Row 5: All X's (X X X X X)  } \\
\\ 
\texttt{Now, please describe your grid.} \\
            }
        }
    }
     \\ \\

    \theutterance \stepcounter{utterance}  
    & & & \multicolumn{2}{p{0.3\linewidth}}{
        \cellcolor[rgb]{0.9,0.9,0.9}{
            \makecell[{{p{\linewidth}}}]{
                \texttt{\tiny{[GM$|$GM]}}
                \texttt{continue} \\
            }
        }
    }
    & & \\ \\

    \theutterance \stepcounter{utterance}  
    & & \multicolumn{4}{p{0.6\linewidth}}{
        \cellcolor[rgb]{0.9,0.9,0.9}{
            \makecell[{{p{\linewidth}}}]{
                \texttt{\tiny{[GM$\rangle$P2]}}
                \texttt{This is my DESCRIPTION: A 5x5 grid where the first row is all blank (▢), the second row has X's in columns 2{-}5, the third row has blanks in columns 1{-}3 and X's in columns 4{-}5, the fourth row has X's in columns 1{-}5, and the fifth row has X's in columns 1{-}5. The pattern is symmetric across the vertical axis, with the central row (third row) having a "U" shape of blanks on the left and X's on the right.} \\
\texttt{Now ask a question in order to find out new aspects of my grid that may be different to your grid. Start with "QUESTION:" and do not add anything else.} \\
            }
        }
    }
    & \\ \\

    \theutterance \stepcounter{utterance}  
    & & & \multicolumn{4}{p{0.6\linewidth}}{
        \cellcolor[rgb]{0.9,0.9,0.9}{
            \makecell[{{p{\linewidth}}}]{
                \texttt{\tiny{[GM$\langle$P2]}}
                \texttt{DESCRIPTION: A 5x5 grid where the first row is all blank (▢), the second row has X's in columns 2{-}5, the third row has blanks in columns 1{-}3 and X's in columns 4{-}5, the fourth row has X's in columns 1{-}5, and the fifth row has X's in columns 1{-}5. The pattern is symmetric across the vertical axis, with the central row (third row) having a "U" shape of blanks on the left and X's on the right.} \\
\\ 
\texttt{QUESTION: Is the first column all blank (▢) in your grid?} \\
            }
        }
    }
     \\ \\

    \theutterance \stepcounter{utterance}  
    & & & \multicolumn{2}{p{0.3\linewidth}}{
        \cellcolor[rgb]{0.9,0.9,0.9}{
            \makecell[{{p{\linewidth}}}]{
                \texttt{\tiny{[GM$|$GM]}}
                \texttt{abort, first word: DESCRIPTION:, but expected QUESTION:} \\
            }
        }
    }
    & & \\ \\

\end{supertabular}
}

\end{document}
